\scanpage{10}
\uline{Ex}\quad\tikz[remember picture,baseline]\node[inner sep=0pt] (afftitle) {Groupe affine \unsure{associé} à un groupe $\mathbb{Z}$};

\diagram{\begin{tikzcd}[ampersand replacement=\&,column sep=small,row sep=1.2em]
\& \& \operatorname{Aff}(\mathbb{Z})\subset\unsure{\mathfrak{S}|\mathcal{S}}_{\mathbb{Z}_{d}}
   \arrow[d,phantom,"{\rotatebox{18}{$=$}}"{description,xshift=-4.5em}] \& \& \\
1 \arrow[r] \& \mathbb{Z} \arrow[r,"i"] \&
\operatorname{Aut}_{\substack{\text{ens à}\\\text{gr d'op}}}
 (\mathbb{Z},\mathbb{Z}_{d}) \arrow[r,"p"] \arrow[d,hook] \&
\operatorname{Aut}(\mathbb{Z}) \arrow[r] \arrow[l,bend left=29,"q"] \& 1\\
\& \& \operatorname{Aut}(\mathbb{Z})\times\mathfrak{S}_{\mathbb{Z}_{d}} \& \&
\end{tikzcd}}

\noindent\begin{minipage}[t]{.50\linewidth}
\[
\begin{aligned}
&\gamma\in\operatorname{Aut}(\mathbb{Z})\\
&q(\gamma)=(\gamma,\gamma)\\
&i(g)=(\mathrm{id}_{\mathbb{Z}},\tau_{g^{\unsure{\prime|.}}}^{\unsure{s|\delta|\deleted{s}}})\,\raisebox{-1.8ex}{$\unsure{.}$}\\
&\boxed{\operatorname{int}(q(\gamma))\mid\mathbb{Z}=\gamma}
\end{aligned}
\]
\end{minipage}\hfill
\begin{minipage}[t]{.45\linewidth}
\diagram{\begin{tikzcd}[ampersand replacement=\&,row sep=.1em]
\underbrace{(\gamma,\rho)\in\operatorname{Aut}(\mathbb{Z},\mathbb{Z}_{d})}\\
\Updownarrow\\
\rho(gx)=\gamma(g)\rho(x)\\
\forall g,x\in\mathbb{Z}
\end{tikzcd}}
\end{minipage}
% The little upper-right mark beside the subscript g and the lower detached
% point after i(g) are retained without assigning them a secure purpose.

\[
\operatorname{Aff}_{\deleted{\unsure{\substack{\text{ens. à}\\\text{gr d'op}}}}}
 (\mathbb{Z})\deleted{\unsure{(\mathbb{Z}_{d})|(G_d)}}\simeq
 \underset{\rotatebox{-15}{$\unsure{-}$}}{\mathbb{Z}}\mathbin{.}\operatorname{Aut}(\mathbb{Z})
 \qquad\text{\unsure{(produit|produit|(produits|produits} $1/2$ direct)}
\]
% Detached short oblique stroke below the first product factor: purpose uncertain.

Si $\Gamma\subset\operatorname{Aut}(\mathbb{Z})$ ($\Gamma$ ss-groupe) on a un ss-groupe naturel de $\deleted{\unsure{\operatorname{Aut}_{\substack{\text{ens}\\\text{gr d'op}}}(G,G_d)}}\operatorname{Aff}(\mathbb{Z})\deleted{\unsure{=|-}}$, savoir $\operatorname{Aff}^{\Gamma}(\mathbb{Z})\simeq\mathbb{Z}\mathbin{.}\Gamma$.

\noindent\tikz[remember picture,baseline]\coordinate (affblocktop);
Hom d'ens à op. avec \unsure{groupes|groupe} \unsure{variables|variable}
\[
(E,\mathbb{G})\underset{(\rho,\gamma)}{\longrightarrow}(E',\unsure{G'|C'})
\]

\deleted{\uline{NB}}

\uline{NB} Si $E'$ est un torseur sous $G'$\inserted{et $E\ne\varnothing$}, alors $\rho$ détermine $\gamma$ car \deleted{\unsure{pour}} si $x\in E$, $\forall g\in G$, $\gamma(g)$ est déterminé en termes de $\rho$ par la condition $\rho(g.x)=\gamma(g).\rho(x)$
\tikz[remember picture,baseline]\coordinate (affblockbottom);
\figure{\begin{tikzpicture}[remember picture,overlay]
\draw ([xshift=-.3em]affblocktop.west) .. controls +(-.7em,0) and +(-.25em,.65em) .. ([xshift=-1em,yshift=-.55em]affblocktop.west)
 -- ([xshift=-1em,yshift=.7em]affblocktop.west |- affblockbottom)
 .. controls +(0,-.55em) and +(-.3em,0) .. ([xshift=-.35em]affblocktop.west |- affblockbottom);
\draw[->] ([xshift=-1em,yshift=-.65em]affblocktop.west)
 .. controls +(-3em,2em) and +(-4em,-1em) .. ([xshift=-2em,yshift=-.45em]afftitle.west)
 -- ([yshift=-.45em]afftitle.west) -- ([yshift=.65em]afftitle.west);
\end{tikzpicture}}

Cas $\mathbb{Z}$ \inserted{\unsure{grpe|gpe|groupe} \unsure{ss-jacent à|s-jacent à|sous-jacent à}} vectoriel $V$ sur corps $R$ (réels p ex)
\[
\begin{aligned}
&\Gamma=\operatorname{Aut}_{R}(V)\subset\operatorname{Aut}_{\mathrm{gr}}(V)\\[.8em]
&\operatorname{Aff}_{R}(V)\simeq V\mathbin{.}\operatorname{Aut}_{R}(V)
\end{aligned}
\]
\[
(\underset{\substack{\rotatebox{90}{$\ni$}\\V}}{z},
 \underset{\substack{\rotatebox{90}{$\ni$}\\\underset{R}{\operatorname{Aut}(V)}}}{\gamma})
 \mathbin{.}(z',\gamma')=(z+\gamma(z'),\gamma\gamma')
\]

Si $V$ muni d'une forme quadratique non \unsure{dég.|dég,}\inserted{$Q$} \unsure{on|$\to$} a le groupe des déplacements (transf. affines qui \tikz[remember picture,baseline]\node[inner sep=0pt] (invword) {\unsure{invarient|invariant}}; la distance)
\figure{\begin{tikzpicture}[remember picture,overlay]
\draw ([xshift=.65em,yshift=-.18em]invword.south west)
 .. controls +(.2em,.06em) and +(-.15em,.03em) .. ([xshift=1.25em,yshift=-.18em]invword.south west);
\end{tikzpicture}}
% Detached short stroke beneath the beginning of the verb, not an asserted underline.
\[
\operatorname{Depl}_{\deleted{\unsure{\text{aff}|\text{Aff}|\text{orth}}}}(V,Q)\simeq V\mathbin{.}\operatorname{Orth}(V,Q)
\]
% archive sheet number: 9
